\section{Exercises}
\label{sec:fol-deduction-05}

\begin{exercise}
\label{ex:fol-deduction-01}
Compute \(\FV(\varphi)\) for
\[
 \varphi=P(x,z)\land\forall x\bigl(Q(x,y)\to\exists yR(x,y,z)\bigr).
\]
Identify the quantifier binding each bound occurrence.
\end{exercise}

\begin{exercise}
\label{ex:fol-deduction-02}
Compute capture-avoiding substitutions
\[
 (\exists yR(x,y))[y/x]
 \quad\text{and}\quad
 (\forall x(P(x)\to Q(y)))[f(x)/y].
\]
Record every bound-variable renaming.
\end{exercise}

\begin{exercise}
\label{ex:fol-deduction-03}
Prove term coincidence and formula coincidence simultaneously for a language
with one binary function and one binary relation, then explain why the same
proof covers arbitrary signatures.
\end{exercise}

\begin{exercise}
\label{ex:fol-deduction-04}
Give a \(\ND_{\FOL}\)-derivation of
\[
 \forall x(P(x)\to Q(x)),\ \exists xP(x)
 \derives_{\ND_{\FOL}}\exists xQ(x).
\]
Mark the eigenvariable and every discharged assumption.
\end{exercise}

\begin{exercise}
\label{ex:fol-deduction-05}
Derive transitivity of equality from \(=I\) and \(=E\).
\end{exercise}

\begin{exercise}
\label{ex:fol-deduction-06}
Translate each sentence using one declared signature: ``Every reviewer read
exactly one paper''; ``One paper was read by every reviewer''; ``No reviewer
read only papers that cited themselves.''  State the domain.
\end{exercise}

\begin{exercise}
\label{ex:fol-deduction-07}
Prove both directions of the claim that the displayed group theory
axiomatizes groups in the signature \(\Lang_{\mathrm{grp}}\).
\end{exercise}

\begin{exercise}
\label{ex:fol-deduction-08}
Write a theory whose models are simple undirected graphs.  Treat equality as
logical identity, and state whether loops are permitted.
\end{exercise}

\begin{exercise}
\label{ex:fol-deduction-09}
Give a cut-free \(\LK_{\FOL}\)-derivation of
\(\forall xP(x)\Rightarrow\exists xP(x)\).  Identify where the nonempty-domain
convention is used.
\end{exercise}

